\section{经典模型：Kim (2014) 的 CNN 文本分类器}

\subsection{模型架构}

Yoon Kim 在 2014 年提出了一个简洁而有效的 CNN 文本分类模型，这是 CNN 用于 NLP 最著名的工作之一。

\begin{figure}[H]
\centering
\includegraphics[width=0.95\textwidth]{slides-images/slide-022.jpg}
\caption{Kim (2014) CNN 模型架构：多种大小的滤波器 + Max Pooling + Softmax 分类\protect\footnotemark}
\end{figure}
\footnotetext{Slides 第23页。图来自 Zhang \& Wallace (2015) 的可视化。}

模型的完整流程如下：

\begin{enumerate}
    \item \textbf{输入层}：句子中每个词用预训练词向量（如 Word2Vec、GloVe）表示
    \item \textbf{卷积层}：使用三种大小的滤波器（bigram、trigram、4-gram），每种 100 个
    \item \textbf{Max Pooling}：对每个滤波器取全局最大值
    \item \textbf{拼接}：将所有滤波器的 max pooling 结果拼接为一个固定长度向量
    \item \textbf{Softmax 分类层}：$y = \text{softmax}(W^{(S)}z + b)$
\end{enumerate}

\begin{importantbox}{模型公式}
对于使用 $m$ 个滤波器、滤波器大小分别为 $h_1, h_2, \ldots$ 的模型：
\begin{enumerate}
    \item 卷积：$\mathbf{c}_j = f(\mathbf{w}_j \cdot \mathbf{x}_{i:i+h-1} + b_j)$，对每个位置 $i$
    \item Max pooling：$\hat{c}_j = \max\{\mathbf{c}_j\}$
    \item 最终特征向量：$\mathbf{z} = [\hat{c}_1, \ldots, \hat{c}_m]$
    \item 分类：$y = \text{softmax}(W^{(S)}\mathbf{z} + b)$
\end{enumerate}
\end{importantbox}

\subsection{多通道词向量：解决微调困境}

Kim 在这篇论文中提出了一个重要的实践技巧：\textbf{双通道词向量}。

\begin{figure}[H]
\centering
\includegraphics[width=0.95\textwidth]{slides-images/slide-023.jpg}
\caption{词向量微调的困境：训练集中出现的词（tedious, dull）向量被更新移动，未出现的词（plodding）留在原位，导致语义关系被破坏\protect\footnotemark}
\end{figure}
\footnotetext{Slides 第24页。}

\begin{warningbox}{词向量微调的陷阱}
当在小型监督数据集上微调预训练词向量时，会出现一个问题：
\begin{itemize}
    \item 训练集中出现的词（如 tedious, dull）的向量被更新，移向``负面情感区域''
    \item 训练集中\textbf{未出现}的同义词（如 plodding）的向量保持不变
    \item 结果是，原本相似的词（tedious 和 plodding）在向量空间中被拉开距离
    \item plodding 可能被错误地分类为正面词
\end{itemize}
这就是为什么微调词向量有时会帮助、有时会损害模型性能。
\end{warningbox}

Kim 的解决方案很简单：使用\textbf{两个通道}的词向量，一个允许微调，一个保持固定（frozen）。每个滤波器同时在两个通道上操作，这样模型可以同时利用：

\begin{itemize}
    \item \textbf{微调通道}：针对具体任务优化的词表示
    \item \textbf{固定通道}：保持原始语义关系的词表示
\end{itemize}

\subsection{实验结果}

\begin{figure}[H]
\centering
\includegraphics[width=0.95\textwidth]{slides-images/slide-024.jpg}
\caption{Kim (2014) 实验结果：在多个文本分类数据集上，简单的 CNN 模型取得了与当时最优方法可比的结果\protect\footnotemark}
\end{figure}
\footnotetext{Slides 第25页。}

Kim 的 CNN 在多个数据集上取得了竞争力强的结果，包括 Stanford Sentiment Treebank（SST）、Movie Reviews（MR）、Subjectivity（Subj）和问题分类（TREC）等。

\begin{figure}[H]
\centering
\includegraphics[width=0.95\textwidth]{slides-images/slide-025.jpg}
\caption{比较需谨慎：Kim (2014) 使用了 Dropout 正则化，而许多对比方法发表于 Dropout 出现之前\protect\footnotemark}
\end{figure}
\footnotetext{Slides 第26页。}

\begin{warningbox}{实验比较中的公平性问题}
Manning 指出了这篇论文实验对比中的一个细微问题：Kim 使用了 Dropout（2012 年提出），而许多对比方法是在 Dropout 出现之前发表的。Dropout 本身可以带来 2--4\% 的准确率提升。更严谨的做法应该是为所有对比方法都加上 Dropout 再进行比较。尽管如此，这项工作仍然证明了简单的 CNN 架构在文本分类任务上的有效性。
\end{warningbox}

\subsection{本章小结}

Kim (2014) 展示了一个极其简洁的 CNN 文本分类架构：多种大小的滤波器 + Max Pooling + Softmax。双通道词向量技巧有效缓解了小数据集上微调词向量的问题。这个模型虽然简单，却在当时取得了接近最优的结果。

%% ========== 模型工具箱比较 ==========

